\section{Conventions and the K3P model}\label{sec:conventions}

\subsection{Rooted and semi-directed networks}

\begin{definition}[Rooted binary network]
Let \(X\) be a finite set of leaf labels.  A rooted binary phylogenetic
network is a finite directed acyclic graph without parallel arcs in which the
root has bidegree \((0,2)\), every tree vertex has bidegree \((1,2)\), every
reticulation has bidegree \((2,1)\), and every leaf has bidegree \((1,0)\) and
a unique label in \(X\).  Every vertex is reachable from the root.  The root
is the lowest stable ancestor of the leaves.  A rooted presentation is
\emph{tree-child} when each internal vertex has a tree vertex or leaf as a
child.
\end{definition}

\begin{definition}[Standard semi-directed convention]\label{def:standard}
Mark every arc entering a reticulation, undirect every other arc, delete the
binary root, and merge its two incident arcs into one edge, retaining at
either endpoint exactly the arrowhead present before deletion.  We admit the
presentation only when this single suppression creates no loop or parallel
edge and loses no reticulation arrowhead.  No later exhaustive cleanup is
part of this operation.  The resulting labelled mixed graph is a
\emph{standard semi-directed network}.  An admissible rooting is a rooted
binary network whose one-step suppression is exactly that mixed graph.
\end{definition}

\begin{figure}[t]
\centering
\begin{tikzpicture}[
  x=1cm,y=1cm,
  every node/.style={font=\scriptsize},
  tvertex/.style={circle,draw,inner sep=1.25pt,minimum size=4.8mm},
  retic/.style={circle,draw,double,inner sep=1.1pt,minimum size=5.2mm},
  leaf/.style={rectangle,draw,rounded corners=.6pt,inner sep=1.5pt},
  rooted/.style={-{Stealth[length=1.9mm,width=1.2mm]},line width=.55pt},
  retained/.style={-{Stealth[length=2.2mm,width=1.35mm]},line width=.9pt},
  ordinary/.style={line width=.65pt}
]
  % Rooted presentation.
  \begin{scope}[shift={(0,0)}]
    \node[font=\small] at (0,3.45) {rooted presentation};
    \node[tvertex] (r) at (0,2.8) {$r$};
    \node[tvertex] (u) at (-1.15,1.75) {$u$};
    \node[tvertex] (v) at (1.15,1.75) {$v$};
    \node[retic] (x) at (0,.65) {$x$};
    \node[leaf] (l0) at (-2.0,.65) {$L_0$};
    \node[leaf] (l1) at (2.0,.65) {$L_1$};
    \node[leaf] (l2) at (0,-.45) {$L_2$};
    \draw[rooted] (r) -- (u);
    \draw[rooted] (r) -- (v);
    \draw[retained] (u) -- (x);
    \draw[retained] (v) -- (x);
    \draw[rooted] (u) -- (l0);
    \draw[rooted] (v) -- (l1);
    \draw[rooted] (x) -- (l2);
  \end{scope}

  \draw[-{Stealth[length=2.3mm,width=1.4mm]},line width=.7pt]
    (2.65,1.4) -- node[above,align=center,font=\tiny]
    {mark reticulation entries\\suppress $r$ once} (4.0,1.4);

  % Standard semi-directed image.
  \begin{scope}[shift={(6.65,0)}]
    \node[font=\small] at (0,3.45) {standard semi-directed graph};
    \node[tvertex] (u2) at (-1.15,2.15) {$u$};
    \node[tvertex] (v2) at (1.15,2.15) {$v$};
    \node[retic] (x2) at (0,.65) {$x$};
    \node[leaf] (l02) at (-2.0,.65) {$L_0$};
    \node[leaf] (l12) at (2.0,.65) {$L_1$};
    \node[leaf] (l22) at (0,-.45) {$L_2$};
    \draw[ordinary] (u2) -- node[above,font=\tiny] {merged root edge} (v2);
    \draw[retained] (u2) -- (x2);
    \draw[retained] (v2) -- (x2);
    \draw[ordinary] (u2) -- (l02);
    \draw[ordinary] (v2) -- (l12);
    \draw[ordinary] (x2) -- (l22);
  \end{scope}
\end{tikzpicture}
\caption{The fixed standard semi-directed convention.  The binary root is
deleted and its two incident arcs are merged in one step; no further cleanup
is performed.  The two arrowheads entering the reticulation are retained.}
\label{fig:standard-semi-directed}
\end{figure}


The one-step operation and the retained reticulation arrowheads are shown in
\cref{fig:standard-semi-directed}.

Isomorphisms preserve leaf labels, ordinary edges, retained arrowheads, and
vertex roles.  Restriction to a leaf subset is a distinct operation: after
choosing displayed reticulation parents one may prune unlabelled dead ends
and suppress unlabelled bivalent vertices.  Such cleanup does not enlarge the
admissible-rooting set of the original mixed graph.  See
\cite{SempleSteel2003,HoltgrefeEtAl2026} for broader background on rooted and
semi-directed conventions.

\begin{definition}[Weak and strong tree-childness]
A standard semi-directed network is \emph{weakly tree-child} if at least one
admissible rooting is tree-child.  It is \emph{strongly tree-child} if it has
an admissible rooting and every admissible rooting is tree-child.  We denote
these classes by \(W_{\mathrm{TC}}\) and \(S_{\mathrm{TC}}\), respectively.
\end{definition}

\begin{lemma}[No-omnian criterion]\label{lem:no-omnian}
In the binary fixed convention, a standard semi-directed network with an
admissible rooting is strongly tree-child if and only if every tail of a
retained reticulation edge is incident with two ordinary edges.
\end{lemma}

\begin{proof}
If a tail fails the condition, it is either a reticulation with a
reticulation child or an ordinary vertex tailing two retained edges.  The
first case is non-tree-child in every rooting.  In the second, moving an
admissible root to the tail's unique ordinary incidence produces an
admissible rooting in which the tail has two reticulation children.  Thus
strong tree-childness fails.  Conversely, under the incidence condition an
ordinary vertex with a reticulation child retains an ordinary child, a
reticulation has an ordinary child, and a root with two reticulation children
would suppress to a forbidden two-arrowhead edge.  Hence every admissible
rooting is tree-child.  This is the fixed binary specialization of the
criterion used in \cite{EnglanderEtAl2026,HoltgrefeEtAl2026}.
\end{proof}

A \emph{bridge} is an edge whose deletion disconnects the underlying graph.
A \emph{blob} is a maximal nontrivial biconnected block.  A network is
\emph{level two} when every blob contains at most two reticulations.
Contracting each nontrivial blob while retaining the intervening ordinary
branching components gives the labelled reduced \emph{tree of blobs}.  When a
blob is cut from the network, each incident bridge becomes a labelled
boundary \emph{port}.  A complete ported factor retains every physical port,
including the child port below a path-sink reticulation.

\begin{definition}[Ordinary triangle equivalence]\label{def:triangle}
An ordinary triangle redirection changes only which vertex of a three-cycle
factor incident with three external arms is reticulate and redirects the two
arrowheads entering it.  Write \(N\trianglerel N'\) when the labelled reduced
trees of blobs agree, corresponding nontriangle complete factors are
labelled mixed-graph isomorphic, and every remaining pair differs by an
ordinary triangle redirection, with coherent boundary transports.
\end{definition}

\subsection{Fourier parameterization and physical domains}

Identify the states and characters with
\(\G=\{0,C,G,T\}\), with Klein-four addition.  A K3P edge \(e\) has Fourier
spectrum

\[
 \widehat f_e=(1,c_e,g_e,t_e).
\]

Inverse Fourier transformation gives

\begin{align}\label{eq:inverse-fourier}
 f_e(0)&=\frac{1+c_e+g_e+t_e}{4},&
 f_e(C)&=\frac{1+c_e-g_e-t_e}{4},\nonumber\\
 f_e(G)&=\frac{1-c_e+g_e-t_e}{4},&
 f_e(T)&=\frac{1-c_e-g_e+t_e}{4}.
\end{align}

The root state is uniform.  At each reticulation \(r\), one parent is chosen
with probability \(\lambda_r\in(0,1)\) and the other with probability
\(1-\lambda_r\).  A switching \(\sigma\) gives a displayed tree \(T_\sigma\)
with positive weight \(w_\sigma\).  For a character assignment
\(\mathbf h=(h_i)_{i\in\X}\),

\begin{equation}\label{eq:k3p-fourier}
 q_N(\mathbf h)=
 \begin{cases}
 \displaystyle\sum_\sigma w_\sigma
  \prod_{e\in T_\sigma}
  \widehat f_e\!\left(\xor_{i\in A_e}h_i\right),
  &\xor_{i\in\X}h_i=0,\\[2mm]
 0,&\text{otherwise},
 \end{cases}
\end{equation}

where \(A_e\mid A_e^c\) is the displayed-tree split.  The map is polynomial
in edge spectra and inheritance probabilities.

Strict positivity of \eqref{eq:inverse-fourier}, together with positive
nontrivial spectra, is exactly the principal domain \(\Dplus\) in
\cref{thm:main}.  Symmetric continuous-time K3P rates
\((\alpha_C,\alpha_G,\alpha_T)>0\) give

\[
 c=e^{-2(\alpha_G+\alpha_T)},\qquad
 g=e^{-2(\alpha_C+\alpha_T)},\qquad
 t=e^{-2(\alpha_C+\alpha_G)},
\]

and are equivalent to the three inequalities defining \(\Dct\).

Write \(\Theta_{3,+}(N)\) and \(\Theta_{3,\mathrm{CT}}(N)\) for the open
physical parameter spaces, \(\Phi_N\) for \eqref{eq:k3p-fourier}, and
\(\M_{3,+}(N)=\Phi_N(\Theta_{3,+}(N))\).  Put

\[
 d_N=\max_{\theta\in\Theta_{3,+}(N)}\rank D\Phi_N(\theta).
\]

A source parameter is regular when its Jacobian has rank \(d_N\).

\subsection{Observational relations}

\begin{definition}[Directed containment and common germ]\label{def:relations}
Write \(N\preceqplus N'\) if there are a regular source point
\(\theta_0\in\Theta_{3,+}(N)\), a connected open neighborhood
\(U\subseteq\Theta_{3,+}(N)\) on which \(\Phi_N\) has rank \(d_N\), and a
real-analytic physical target section
\(\sigma:U\to\Theta_{3,+}(N')\) such that

\[
 \Phi_N=\Phi_{N'}\circ\sigma\qquad\text{on }U.
\]

Write \(N\bowtieplus N'\) when the two physical images contain a common
analytic germ that is regular and full-dimensional in both images and has
physical analytic sections from both parameter spaces.  Define
\(\preceqct\) and \(\bowtiect\) by replacing the edge domain with \(\Dct\).
\end{definition}

These relations are stronger than nonempty image intersection and different
from inclusion of complex Zariski closures.  The target section in
\(\preceqplus\) is physical but need not be target-open or target-regular.
This asymmetric quantifier is important in the cut-transfer argument below.
